%
% teil2.tex -- Beispiel-File für teil2 
%
% (c) 2020 Prof Dr Andreas Müller, Hochschule Rapperswil
%
% !TEX root = ../../buch.tex
% !TEX encoding = UTF-8
%
\subsection{Mathematische Vertiefung: Die komplexe Konturintegration
    \label{kepler:subsection:komplexe_integration}}

Bei der rigorosen analytischen Berechnung der Umlaufzeit für eine elliptische Bahn stösst man zwangsläufig auf ein spezifisches trigonometrisches Integral \cite{nahin2010}.
Dieses Integral für Umlaufbahnen mit der Exzentrizität \(E\) lautet
\begin{equation}
    \int_{0}^{2\pi} \frac{d\theta}{(1+E \cos(\theta))^2} = \frac{2\pi}{(1-E^2)^{3/2}} .
    \label{kepler:equation:umlauf_integral}
\end{equation}
Um dieses zunächst anspruchsvolle reelle Integral zu lösen, bietet sich ein eleganter Weg über die Transformation in die komplexe Ebene an \cite{nahin2010}.
Man substituiert die Integrationsvariable mit Hilfe der eulerschen Identität durch eine komplexe Zahl \(z = e^{i\theta}\), woraus direkt \(dz = iz d\theta\) folgt \cite{nahin2010}.
Die reelle Cosinus-Funktion lässt sich daraufhin algebraisch als
\begin{equation}
    \cos(\theta) = \frac{z + \frac{1}{z}}{2}
    \label{kepler:equation:cosinus_komplex}
\end{equation}
ausdrücken \cite{nahin2010}.
Durch diese Substitution wird das ursprüngliche reelle Integral zu einem Konturintegral entlang des Einheitskreises in der komplexen Ebene \cite{nahin2010}.
Der umgeformte komplexe Integrand besitzt dabei Singularitäten beziehungsweise Polstellen zweiter Ordnung \cite{nahin2010}.
Mit Hilfe der \emph{cauchyschen Integralformel} für Ableitungen lässt sich das Integral schliesslich exakt auswerten und liefert exakt den Ausdruck auf der rechten Seite von Gleichung \eqref{kepler:equation:umlauf_integral} \cite{nahin2010}.

